\documentclass[10pt, aspectratio=169]{beamer}
% Preámbulo modular para presentaciones Beamer (Modelación Estadística)
% Diseño y paleta institucional del Tecnológico de Monterrey

\documentclass[aspectratio=169,xcolor={dvipsnames,table}]{beamer}

\usepackage[utf8]{inputenc}
\usepackage[spanish,mexico]{babel}
\usepackage{amsmath,amssymb,amsfonts,mathtools}
\usepackage{graphicx}
\usepackage{listings}
\usepackage{booktabs}
\usepackage{tikz}
\usetikzlibrary{matrix,arrows,decorations.pathmorphing,shapes.geometric,positioning}

% Paleta Institucional Tecnológico de Monterrey
\definecolor{TecAzul}{HTML}{0039A6}       % Azul TEC: color primario institucional
\definecolor{TecAzulOscuro}{HTML}{1A2E51} % Azul marino oscuro
\definecolor{TecRojo}{HTML}{EC2661}       % Rojo de acento (paleta miTec)
\definecolor{TecGris}{HTML}{646464}       % Gris neutro para comentarios y subtítulos
\definecolor{TecFondoBloque}{HTML}{F4F6F9}% Fondo suave para bloques

% Configuración de Tema Beamer
\usetheme{Boadilla}
\usecolortheme[named=TecAzulOscuro]{structure}
\setbeamercolor{alerted text}{fg=TecRojo}
\setbeamercolor{palette primary}{bg=TecAzulOscuro,fg=white}
\setbeamercolor{palette secondary}{bg=TecAzul,fg=white}
\setbeamercolor{palette tertiary}{bg=TecAzulOscuro!80!black,fg=white}
\setbeamercolor{title}{fg=TecAzulOscuro,bg=TecFondoBloque}
\setbeamercolor{frametitle}{fg=TecAzulOscuro,bg=TecFondoBloque}

% Bloques personalizados elegantes
\setbeamercolor{block title}{bg=TecAzulOscuro,fg=white}
\setbeamercolor{block body}{bg=TecFondoBloque,fg=black}
\setbeamercolor{block title alerted}{bg=TecRojo,fg=white}
\setbeamercolor{block body alerted}{bg=TecRojo!8,fg=black}
\setbeamercolor{block title example}{bg=TecAzul,fg=white}
\setbeamercolor{block body example}{bg=TecAzul!8,fg=black}

% Redondear bordes y añadir sombra sutil
\setbeamertemplate{blocks}[rounded][shadow=true]
\setbeamertemplate{navigation symbols}{}

% Pie de página limpio con número de diapositiva
\setbeamertemplate{footline}{
  \leavevmode%
  \hbox{%
  \begin{beamercolorbox}[wd=.7\paperwidth,ht=2.25ex,dp=1ex,left]{author in head/foot}%
    \hspace*{2ex}\insertshortauthor~~(\insertshortinstitute) \hfill \insertshorttitle
  \end{beamercolorbox}%
  \begin{beamercolorbox}[wd=.3\paperwidth,ht=2.25ex,dp=1ex,right]{date in head/foot}%
    \insertshortdate{}\hspace*{2em}
    \insertframenumber{} / \inserttotalframenumber\hspace*{2ex}
  \end{beamercolorbox}}%
  \vskip0pt%
}

% Configuración de Listings de Python para visualización pedagógica de código
\lstset{
    language=Python,
    basicstyle=\ttfamily\footnotesize,
    keywordstyle=\color{TecAzulOscuro}\bfseries,
    stringstyle=\color{TecRojo},
    commentstyle=\color{TecGris}\itshape,
    showstringspaces=false,
    breaklines=true,
    frame=lines,
    rulecolor=\color{TecAzulOscuro!30},
    backgroundcolor=\color{TecFondoBloque},
    aboveskip=0.5em,
    belowskip=0.5em,
    literate={á}{{\'a}}1 {é}{{\'e}}1 {í}{{\'i}}1 {ó}{{\'o}}1 {ú}{{\'u}}1
             {Á}{{\'A}}1 {É}{{\'E}}1 {Í}{{\'I}}1 {Ó}{{\'O}}1 {Ú}{{\'U}}1
             {ñ}{{\~n}}1 {Ñ}{{\~N}}1 {¿}{{?`}}1 {¡}{{!`}}1
}

% Metadatos institucionales por defecto
\title[Modelación Estadística]{Modelación Estadística para Ciencia de Datos e Ingeniería}
\author[J. Castillo Colmenares]{Juliho David Castillo Colmenares}
\institute[Tec de Monterrey]{Tecnológico de Monterrey \\ Escuela de Ingeniería y Ciencias}
\date{\today}

% Comandos y macros estadísticas compartidas para las presentaciones Beamer (Metropolis)

% Conjuntos numéricos básicos
\newcommand{\R}{\mathbb{R}}
\newcommand{\N}{\mathbb{N}}
\newcommand{\Q}{\mathbb{Q}}
\newcommand{\Z}{\mathbb{Z}}
\newcommand{\set}[1]{\left\{ #1 \right\}}
\newcommand{\abs}[1]{\left|#1\right|}
\newcommand{\norm}[1]{\left\| #1 \right\|}

% Comandos estadísticos y probabilísticos del curso
\newcommand{\Var}{\operatorname{Var}}
\newcommand{\cov}{\operatorname{Cov}}
\newcommand{\corr}{\rho}
\newcommand{\comb}[2]{\begin{pmatrix} #1 \\ #2 \end{pmatrix}}
\newcommand{\s}{\sigma}
\newcommand{\particion}{\mathfrak{P}}
\newcommand{\card}[1]{\#\left( #1 \right)}
\newcommand{\seq}[3]{\set{#1}_{#2}^{#3}}
\newcommand{\E}{\mathbb{E}}
\newcommand{\Prob}{\mathbb{P}}

% Entornos pedagógicos y cajas en español académico natural
\newenvironment{discusion}{%
  \begin{alertblock}{Pregunta de Discusión}%
}{%
  \end{alertblock}%
}

\newenvironment{reflexion}{%
  \begin{alertblock}{Reflexión para el Aula}%
}{%
  \end{alertblock}%
}

\newenvironment{paradoja}{%
  \begin{alertblock}{Paradoja de Exploración}%
}{%
  \end{alertblock}%
}

\newenvironment{motivacion}{%
  \begin{exampleblock}{Motivación: Ciencia de Datos en la Práctica}%
}{%
  \end{exampleblock}%
}

\newenvironment{retoclase}{%
  \begin{block}{Reto Rápido en Equipo}%
}{%
  \end{block}%
}


\title[Regresión Lineal Simple]{Sección 09.01: Regresión Lineal Simple (MCO)}
\subtitle{Mínimos Cuadrados Ordinarios y el Coeficiente de Determinación $R^2$}
\author[J. Castillo Colmenares]{Juliho Castillo Colmenares}
\institute{Tecnológico de Monterrey}
\date{\vspace{-1.2cm}}

\begin{document}

% Slide 1: Portada
\begin{frame}[plain]
  \titlepage
\end{frame}

% Slide 2: Hoja de Ruta
\begin{frame}{Hoja de Ruta --- Capítulo 09: Regresiones Lineales y Múltiples}
  \footnotesize
  ¿Dónde nos encontramos en el desarrollo modular del capítulo? \pause
  \begin{itemize}\itemsep=0.08em
    \item \textbf{09.01} Regresión Lineal Simple (MCO) y Coeficiente $R^2$ \emph{(Hoy)} \pause
    \item \textbf{09.02} Regresión Lineal Múltiple, Ecuación Normal y Ridge/Lasso \pause
    \item \textbf{09.03} Diagnóstico de Residuos, Multicolinealidad (VIF) y Supuestos \pause
    \item \textbf{09.04} Validación de Modelos, $k$-fold Cross-Validation y \texttt{scikit-learn}
  \end{itemize}
  \vspace{-0.05cm}
  \begin{block}{Objetivo de la Sesión}
    Generalizar la correlación (Capítulo 03) a un modelo predictivo: derivar los estimadores de Mínimos Cuadrados Ordinarios (MCO) para la recta $\hat Y=\hat\beta_0+\hat\beta_1 X$, descomponer la variabilidad total en variabilidad explicada y residual, y evaluar la significancia estadística de la pendiente.
  \end{block}
\end{frame}

% Slide 3: Motivación
\begin{frame}{De la Correlación a la Regresión}
  \footnotesize
  \begin{motivacion}
    La correlación de Pearson $r_{XY}$ mide la \emph{fuerza y dirección} de una asociación lineal entre $X$ y $Y$, pero no permite \emph{predecir} valores nuevos. La regresión lineal responde una pregunta distinta: dado un valor de $X$, ¿cuál es la mejor estimación puntual de $Y$?
  \end{motivacion}

  \vspace{0.15cm}
  \pause
  \begin{itemize}\itemsep=0.1cm
    \item El modelo propuesto es $Y_e=\hat\beta_0+\hat\beta_1 X$: una recta que resume la tendencia central de la nube de puntos. \pause
    \item El criterio de ajuste es minimizar la suma de cuadrados de los residuos $\sum(Y-Y_e)^2$ (Mínimos Cuadrados Ordinarios).
  \end{itemize}
\end{frame}

% Slide 4: El modelo y el objetivo
\begin{frame}{El Modelo Lineal y la Función Objetivo}
  \footnotesize
  \begin{block}{Modelo Poblacional}
    $$Y_i=\beta_0+\beta_1 X_i+\varepsilon_i, \qquad E[\varepsilon_i]=0,\ \text{Var}(\varepsilon_i)=\sigma^2,\ \text{Cov}(\varepsilon_i,\varepsilon_j)=0$$
  \end{block}
  \pause

  \vspace{0.1cm}
  \begin{alertblock}{Función Objetivo de Mínimos Cuadrados}
    $$Q(\beta_0,\beta_1)=\sum_{i=1}^n\left(Y_i-(\beta_0+\beta_1 X_i)\right)^2 \quad \longrightarrow \quad \min_{\beta_0,\beta_1} Q$$
  \end{alertblock}
\end{frame}

% Slide 5: Derivación de las fórmulas
\begin{frame}{Derivación de las Ecuaciones Normales de Gauss}
  \footnotesize
  Igualando a cero las derivadas parciales $\partial Q/\partial\beta_0=0$ y $\partial Q/\partial\beta_1=0$: \pause

  \vspace{0.1cm}
  \begin{block}{Estimadores de Mínimos Cuadrados}
    $$\hat\beta_1=\frac{\sum_{i=1}^n(X_i-\bar X)(Y_i-\bar Y)}{\sum_{i=1}^n(X_i-\bar X)^2}=\frac{S_{XY}}{S_{XX}}, \qquad \hat\beta_0=\bar Y-\hat\beta_1\bar X$$
  \end{block}
  \pause

  \vspace{0.05cm}
  \begin{alertblock}{Propiedades de Proyección Ortogonal}
    Los residuos $\hat\varepsilon_i=Y_i-\hat Y_i$ satisfacen exactamente $\sum\hat\varepsilon_i=0$ y $\sum X_i\hat\varepsilon_i=0$: la recta ajustada siempre pasa por $(\bar X,\bar Y)$.
  \end{alertblock}
\end{frame}

% Slide 6: Descomposición y R²
\begin{frame}{Descomposición de la Varianza y el Coeficiente $R^2$}
  \footnotesize
  \begin{block}{Teorema de Partición}
    $$\underbrace{\sum(Y_i-\bar Y)^2}_{\text{SCT}} = \underbrace{\sum(\hat Y_i-\bar Y)^2}_{\text{SCR}} + \underbrace{\sum(Y_i-\hat Y_i)^2}_{\text{SCE}}$$
  \end{block}
  \pause

  \vspace{0.1cm}
  \begin{alertblock}{Coeficiente de Determinación}
    $$R^2=\frac{\text{SCR}}{\text{SCT}}=1-\frac{\text{SCE}}{\text{SCT}} \in[0,1], \qquad \text{en regresión simple: } R^2=r_{XY}^2$$
  \end{alertblock}
\end{frame}

% Slide 7: Significancia de la pendiente
\begin{frame}{Significancia Estadística de la Pendiente}
  \footnotesize
  \begin{block}{Prueba de Hipótesis}
    $$H_0:\beta_1=0 \quad (\text{sin relación lineal}) \qquad \text{vs.} \qquad H_a:\beta_1\neq0$$
  \end{block}
  \pause

  \vspace{0.1cm}
  \begin{alertblock}{Estadístico de Prueba}
    $$t=\frac{\hat\beta_1}{SE(\hat\beta_1)}, \qquad SE(\hat\beta_1)=\sqrt{\frac{\text{CME}}{S_{XX}}} \ \sim\ t_{n-2} \quad\text{bajo } H_0$$
  \end{alertblock}
\end{frame}

% Slide 8: Lab Python (Bloque 1)
\begin{frame}[fragile]{Laboratorio en Python: Estimación MCO}
  \scriptsize
  Estimación de $\hat\beta_0,\hat\beta_1$ verificada contra \texttt{scipy.stats.linregress} (\texttt{09.01\_simple\_linear\_regression.py}):
  {\lstset{basicstyle=\fontsize{3.6pt}{4.3pt}\selectfont\ttfamily,aboveskip=0.05em,belowskip=0.05em}
  \lstinputlisting[firstline=17, lastline=32]{../../code/09_regresiones/09.01_simple_linear_regression.py}}
\end{frame}

% Slide 9: Lab Python (Bloque 2)
\begin{frame}[fragile]{Laboratorio en Python: Descomposición y $R^2$}
  \scriptsize
  Verificación de SCT=SCR+SCE y de la identidad $R^2=r^2$:
  {\lstset{basicstyle=\fontsize{3.8pt}{4.5pt}\selectfont\ttfamily,aboveskip=0.05em,belowskip=0.05em}
  \lstinputlisting[firstline=35, lastline=48]{../../code/09_regresiones/09.01_simple_linear_regression.py}}
\end{frame}

% Slide 10: Lab Python (Bloque 3)
\begin{frame}[fragile]{Laboratorio en Python: Prueba $t$ e Intervalos}
  \scriptsize
  Prueba de significancia de la pendiente e intervalos de confianza/predicción:
  {\lstset{basicstyle=\fontsize{3.4pt}{4.1pt}\selectfont\ttfamily,aboveskip=0.05em,belowskip=0.05em}
  \lstinputlisting[firstline=51, lastline=71]{../../code/09_regresiones/09.01_simple_linear_regression.py}}
\end{frame}

% Slide 11: Salida Terminal
\begin{frame}[fragile]{Laboratorio en Python: Salida en Terminal}
  \scriptsize
  Salida estándar generada al ejecutar el script del laboratorio en Python:
  \vspace{0.02cm}
  \begin{block}{Salida Estándar en Consola}
    \fontsize{3.8pt}{4.6pt}\selectfont
    \begin{verbatim}
=== Block 1: OLS Coefficient Estimation ===
beta_1_hat=3.0910, beta_0_hat=6.8202 (matches scipy.stats.linregress)

=== Block 2: Sum-of-Squares Decomposition ===
SST=4592.92 = SSR(4422.38)+SSE(170.54); R^2=0.9629 = r^2 (Pearson r=0.9813)

=== Block 3: Slope Significance and Intervals ===
t=18.3608 (df=13), p=0.000000
95% CI mean: (41.86, 45.97);  95% PI new obs: (35.82, 52.00)
    \end{verbatim}
  \end{block}
\end{frame}

% Slide 12A: Ejercicio Nivel 1 (Enunciado)
\begin{frame}{Ejercicio en Clase --- Nivel 1: Fundamental (1/2)}
  \footnotesize
  \begin{block}{Problema (Enunciado, Problema 9.8.1)}
    Se estudia la relación entre horas de estudio ($X$) y calificación ($Y$) en $n=10$ estudiantes: $X=\{2,3,4,5,6,7,8,9,10,12\}$, $Y=\{45,50,55,60,65,70,75,80,85,90\}$. Calcule $\hat\beta_0,\hat\beta_1$, prediga la calificación para $X=7.5$ horas, y calcule $R^2$.
  \end{block}

  \vspace{0.15cm}
  \pause
  \begin{alertblock}{Planteamiento}
    Calcule $\bar X,\bar Y,S_{XX},S_{XY}$ a partir de las 10 parejas de datos.
  \end{alertblock}
\end{frame}

% Slide 12B: Ejercicio Nivel 1 (Resolución)
\begin{frame}{Ejercicio en Clase --- Nivel 1: Fundamental (2/2)}
  \scriptsize
  Los datos son casi perfectamente lineales: \pause

  \vspace{0.05cm}
  \begin{block}{Cálculo}
    $$\bar X=6.6,\ \bar Y=67.5, \qquad \hat\beta_1\approx4.708,\ \hat\beta_0\approx36.43 \implies \hat Y=36.43+4.708X$$
    $$\hat Y(7.5)\approx71.74, \qquad R^2\approx0.993$$
  \end{block}

  \vspace{0.05cm}
  \pause
  \begin{alertblock}{Interpretación}
    Un estudiante con 7.5 horas de estudio obtendría una calificación esperada de $\approx71.7$. El $R^2\approx0.993$ indica que el modelo lineal explica prácticamente toda la variabilidad de las calificaciones.
  \end{alertblock}
\end{frame}

% Slide 13A: Ejercicio Nivel 2 (Enunciado)
\begin{frame}{Ejercicio en Clase --- Nivel 2: Operativo (1/2)}
  \footnotesize
  \begin{block}{Problema (Enunciado, Problema 9.8.5)}
    Con $n=15$ días, $\bar X=25.0°C$, $\bar Y=450.0$ MWh, $S_{XX}=500.0$, $S_{XY}=2000.0$ y $\text{SCT}=10000.0$. Calcule $\hat\beta_0,\hat\beta_1$, la $\text{SCR}$ y $\text{SCE}$, y pruebe $H_0:\beta_1=0$ al $\alpha=0.05$ ($t_{0.025,13}=2.160$).
  \end{block}

  \vspace{0.15cm}
  \pause
  \begin{alertblock}{Estrategia}
    \scriptsize
    $\hat\beta_1=S_{XY}/S_{XX}$; $\text{SCR}=\hat\beta_1 S_{XY}$; $\text{SCE}=\text{SCT}-\text{SCR}$.
  \end{alertblock}
\end{frame}

% Slide 13B: Ejercicio Nivel 2 (Resolución)
\begin{frame}{Ejercicio en Clase --- Nivel 2: Operativo (2/2)}
  \scriptsize
  Calculamos la pendiente, la partición de varianza y el estadístico $t$: \pause

  \vspace{0.05cm}
  \begin{block}{Cálculo}
    $$\hat\beta_1=\frac{2000.0}{500.0}=4.0, \quad \hat\beta_0=450.0-4.0(25.0)=350.0$$
    $$\text{SCR}=4.0(2000.0)=8000.0, \quad \text{SCE}=10000.0-8000.0=2000.0, \quad \text{CME}=\frac{2000.0}{13}\approx153.85$$
    $$SE(\hat\beta_1)=\sqrt{\frac{153.85}{500.0}}\approx0.5547, \qquad t=\frac{4.0}{0.5547}\approx7.211$$
  \end{block}

  \vspace{0.05cm}
  \pause
  \begin{alertblock}{Interpretación}
    \scriptsize
    Como $7.215>2.160$, \textbf{se rechaza $H_0$}: la temperatura ambiente tiene un efecto lineal estadísticamente significativo sobre el consumo eléctrico.
  \end{alertblock}
\end{frame}

% Slide 14A: Ejercicio Nivel 3 (Enunciado)
\begin{frame}{Ejercicio en Clase --- Nivel 3: Analítico (1/2)}
  \footnotesize
  \begin{block}{Problema --- Partición y $R^2=r^2$ (Enunciado, Problema 9.8.7)}
    Con $n=30$, $S_X=4.5$, $S_Y=12.0$, $\bar X=8.0\%$, $\bar Y=15.0\%$ y $r_{XY}=0.820$. Demuestre que $R^2=r_{XY}^2$, calcule $\text{SCT}$, $\text{SCR}$, $\text{SCE}$, y obtenga $\hat\beta_1=r_{XY}(S_Y/S_X)$.
  \end{block}

  \vspace{0.1cm}
  \pause
  \begin{alertblock}{Estrategia}
    \scriptsize
    $\text{SCT}=(n-1)S_Y^2$; use la relación algebraica entre $r_{XY}$ y $\hat\beta_1$ a través de las desviaciones estándar.
  \end{alertblock}
\end{frame}

% Slide 14B: Ejercicio Nivel 3 (Resolución)
\begin{frame}{Ejercicio en Clase --- Nivel 3: Analítico (2/2)}
  \scriptsize
  Calculamos la partición de varianza y verificamos la identidad: \pause

  \vspace{0.05cm}
  \begin{block}{Cálculo}
    \scriptsize
    $$\text{SCT}=(30-1)(12.0)^2=4176.0, \qquad R^2=r_{XY}^2=(0.820)^2=0.6724$$
    $$\text{SCR}=R^2\cdot\text{SCT}\approx2807.9, \qquad \text{SCE}=\text{SCT}-\text{SCR}\approx1368.1$$
    $$\hat\beta_1=0.820\left(\frac{12.0}{4.5}\right)\approx2.187, \qquad \hat\beta_0=15.0-2.187(8.0)\approx-2.493$$
  \end{block}

  \vspace{0.05cm}
  \pause
  \begin{alertblock}{Interpretación}
    \scriptsize
    La identidad $R^2=r_{XY}^2$ es exacta en regresión simple: ambos coeficientes cuantifican la misma proporción de variabilidad compartida, solo que $r$ conserva el signo de la asociación mientras $R^2$ no.
  \end{alertblock}
\end{frame}

% Slide 15A: Ejercicio Nivel 4 (Enunciado)
\begin{frame}{Ejercicio en Clase --- Nivel 4: Desafiante (1/2)}
  \footnotesize
  \begin{block}{Problema --- Ecuaciones Normales de Gauss e Insesgo (Enunciado, Problema 9.8.9)}
    Bajo el modelo poblacional $Y_i=\beta_0+\beta_1 X_i+\varepsilon_i$ con $E[\varepsilon_i]=0$, deduzca el sistema de ecuaciones normales de Gauss a partir de $\partial Q/\partial\beta_0=0$ y $\partial Q/\partial\beta_1=0$, y demuestre que $E[\hat\beta_1]=\beta_1$ y $E[\hat\beta_0]=\beta_0$.
  \end{block}

  \vspace{0.1cm}
  \pause
  \begin{alertblock}{Estrategia}
    \scriptsize
    Sustituya el modelo verdadero $Y_i=\beta_0+\beta_1X_i+\varepsilon_i$ dentro de la fórmula cerrada de $\hat\beta_1$ y aplique linealidad de $E[\cdot]$.
  \end{alertblock}
\end{frame}

% Slide 15B: Ejercicio Nivel 4 (Resolución)
\begin{frame}{Ejercicio en Clase --- Nivel 4: Desafiante (2/2)}
  \scriptsize
  Sustituimos el modelo verdadero en el estimador de pendiente: \pause

  \vspace{0.05cm}
  \begin{block}{Cálculo}
    \scriptsize
    $$\hat\beta_1=\sum_i c_i Y_i, \quad c_i=\frac{X_i-\bar X}{S_{XX}} \implies \hat\beta_1=\sum_i c_i(\beta_0+\beta_1X_i+\varepsilon_i)=\beta_0\underbrace{\sum c_i}_{0}+\beta_1\underbrace{\sum c_iX_i}_{1}+\sum c_i\varepsilon_i$$
    $$E[\hat\beta_1]=\beta_1+\sum c_i E[\varepsilon_i]=\beta_1+0=\beta_1$$
  \end{block}

  \vspace{0.05cm}
  \pause
  \begin{alertblock}{Interpretación}
    \scriptsize
    Como $\hat\beta_0=\bar Y-\hat\beta_1\bar X$ y $E[\bar Y]=\beta_0+\beta_1\bar X$, se sigue directamente $E[\hat\beta_0]=\beta_0$. Ambos estimadores MCO son insesgados sin necesidad de asumir normalidad de los errores.
  \end{alertblock}
\end{frame}

% Slide 16: Cierre
\begin{frame}{Cierre de Sección: Regresión Lineal Simple}
  \begin{reflexion}
    Hemos generalizado la correlación a un modelo predictivo: los estimadores de Mínimos Cuadrados Ordinarios minimizan la suma de cuadrados de los residuos, la descomposición $\text{SCT}=\text{SCR}+\text{SCE}$ cuantifica cuánta variabilidad explica el modelo ($R^2$), y la prueba $t$ certifica si esa relación es estadísticamente significativa.
  \end{reflexion}

  \vspace{0.2cm}
  \pause
  \begin{block}{Lecciones Clave}
    \begin{itemize}\itemsep=0.08cm
      \item \textbf{MCO:} $\hat\beta_1=S_{XY}/S_{XX}$, $\hat\beta_0=\bar Y-\hat\beta_1\bar X$, ambos insesgados.
      \item \textbf{$R^2=r^2$:} en regresión simple, el coeficiente de determinación es el cuadrado de la correlación de Pearson.
      \item \textbf{Significancia:} $H_0:\beta_1=0$ se contrasta con $t=\hat\beta_1/SE(\hat\beta_1)\sim t_{n-2}$.
    \end{itemize}
  \end{block}
\end{frame}

% Slide 17: Perspectiva Modular
\begin{frame}{Perspectiva Modular --- El Próximo Reto}
  \scriptsize
  \begin{retoclase}
    Con la regresión de un solo predictor dominada, la Sección 09.02 generaliza el modelo a \textbf{múltiples} variables predictoras simultáneas mediante la Ecuación Normal en notación matricial $\hat{\boldsymbol\beta}=(X^TX)^{-1}X^T\mathbf{Y}$, e introduce la regularización Ridge y Lasso para controlar la complejidad del modelo cuando hay muchos predictores correlacionados.
  \end{retoclase}

  \vspace{0.08cm}
  \pause
  \begin{alertblock}{Preparación Recomendada}
    \begin{itemize}\itemsep=0.06cm
      \item Repasa la notación de álgebra lineal (matrices, transpuesta, inversa) que sustituirá a las sumatorias escalares.
      \item Reflexiona sobre por qué $R^2$ nunca disminuye al agregar predictores, y qué problema resuelve el $R^2$ ajustado.
    \end{itemize}
  \end{alertblock}
\end{frame}

\end{document}
